\section{Introduction}

Vision-language models are trained until a picture and a sentence can be
reasoned about together, and are then described as having a shared
representation space. It is not obvious what that claim asserts. A weak reading
is that the two streams are mutually legible: the model can answer questions
about an image in text. That reading is clearly true and equally clearly
uninteresting, because a captioning pipeline with a frozen encoder satisfies it.
A strong reading is that the model's internal state does not depend on which
modality delivered a piece of content --- that content is represented, and the
channel it arrived through is not. This paper is about measuring the distance
between those two readings.

\paragraph{Equivalence is ill-posed as usually stated.}
The natural experiment --- express the same thing in text and in pixels, compare
the representations --- does not survive contact with the details. A noun phrase
and the image region it refers to are not the same object: the phrase
\emph{a dog} denotes a class, the crop denotes a particular animal in a
particular pose. The two occupy different numbers of tokens, so there is no
positional correspondence to compare along. And any residual distance mixes
together at least three things: divergence in how the model encodes the two
channels, genuine information asymmetry between a word and a picture, and
whatever variation the model would show between two \emph{text} expressions of
the same content. A number that confounds these is not evidence about
representation. Much of the reported ``modality gap'' literature measures raw
distances of this kind.

We take the position that the confounds must be designed out rather than
argued away, and that doing so is the contribution. Three constructions do the
work.

\paragraph{Assert equivalence at the merge position, not on the span.}
Two views of an item share a prefix and a suffix verbatim and differ only in a
substituted span. Because the span occupies a different number of tokens in each
view, we make no claim about the span's own representation. We instead compare
the hidden state at the \emph{first token after} the span, located as
$\text{len}(\text{sequence}) - \text{len}(\text{suffix tokens})$. If the model
has reconciled the two views, it has done so by that point; if it has not, the
divergence is still present there, since everything downstream is conditioned on
it. This makes the comparison well-defined without requiring the two views to
have the same length, and it is model-agnostic: it never needs to know how the
processor expanded the image.

\paragraph{Normalise by a within-modality control.}
An absolute distance is uninterpretable because we have no scale for it. We
therefore measure, on the same items, how far apart two \emph{same-modality}
expressions of the same content sit --- a synonym in the text view, a re-render
in the image view --- and report the ratio. The resulting Modality Substitution
Gap answers a question with a meaningful zero point: \emph{does crossing
modalities cost more than re-expressing content within one?} A value near $1$
says it does not. This is the quantity a training objective could sensibly
target, and the denominator is what makes it a claim rather than a number.

\paragraph{Make the substitution exactly content-preserving.}
Referential grounding is an unreliable probe precisely because the phrase and
the region genuinely differ in content, so a nonzero gap is expected under any
hypothesis. Our centrepiece instead substitutes a word with \emph{an image of
that word}, rendered as glyphs. The content is then identical by construction ---
the same lexical item, delivered through the vision encoder --- and the
substitution has no information asymmetry left to explain a gap. It also removes
a limitation that confines grounding-based probes to depictable nouns: an
abstract word has no region to point at, but it renders as readily as a concrete
one, so the same instrument covers function words, abstract nouns, and rare
words. We complement this with a referential tier (as a control, not a headline)
and a relational tier.

\paragraph{The controls are the difficult part.}
An instrument that reports a confident number about an unreadable stimulus is
worse than no instrument. Two controls turn out to be load-bearing, and neither
is standard practice.

\emph{A span-free floor.} If a comprehension probe is used to establish that the
model recovers the substituted content, one must first establish that the probe
depends on that content at all. We score the same forced choice with the span
blanked in both modalities. Chance is the wrong reference here --- distractors
drawn from the same class in a constraining context are not indistinguishable
without the span --- and in our own development this control overturned a
diagnosis we had already committed to.

\emph{A hierarchy of nulls.} A gap in representation is not automatically a gap
in information. If the two modalities' states differ by a constant offset, a
rotation, or any invertible linear map, then no information is lost and any
downstream linear readout can undo the difference; indeed the projector in a
standard VLM \emph{is} such a map. We therefore report $\msg$ after removing, in
turn, a per-modality translation, an orthogonal map, and an arbitrary linear map,
each fitted out-of-fold. The level at which a gap disappears is the result. We
report this hierarchy because it is the analysis most likely to deflate our own
headline number, and because a fitted map with more parameters than fitting data
fails for reasons of sample size in a way that mimics an irreducible gap.

\paragraph{Contributions.}
\begin{enumerate}
    \item A well-posed formulation of cross-modal token substitutability: the
    merge position (\S\ref{sec:merge}) and the normalised Modality Substitution
    Gap (\S\ref{sec:msg}).
    \item Orthographic substitution (\S\ref{sec:tiers}), a content-preserving
    cross-modal probe that extends to words with no visual referent, together
    with referential and relational tiers.
    \item The validity apparatus: a span-free floor for comprehension probes,
    and an out-of-fold null hierarchy separating a gap in information from a gap
    in coordinates (\S\ref{sec:validity}).
    \item A measurement of all of the above on a 2B-parameter VLM, with
    pre-registered decision thresholds fixed before the numbers existed.
\end{enumerate}
